\famichapter{Requirements and Scope}{This chapter specifies the functional and
non-functional requirements of FamiPet, the user roles and their permissions, the
scope of the system, and the use cases that the implementation supports.}{ch:req}

\section{Actors and Roles}

The system distinguishes three actors, defined by the actual implementation.

\begin{table}[htbp]
\centering
\begin{tabular}{@{}p{2.5cm}p{3.2cm}p{8.6cm}@{}}
\toprule
\textbf{Role} & \textbf{Stored as} & \textbf{Permissions enforced by the backend} \\
\midrule
Visitor & no account & Browse public content: pets, breeds, veterinarian
directory, community posts, and lost-and-found reports. Cannot create or modify
any resource. \\
\addlinespace
Registered user & \texttt{User.role = "user"} & All visitor permissions plus
account management, pet management with QR identity, health records,
vaccinations, appointments, reminders, favorites, community posts, lost-and-found
reports, adoption applications, notifications, and the AI assistant. \\
\addlinespace
Administrator & \texttt{User.role = "admin"} & All user permissions plus the
admin panel: dashboard statistics, user management and blocking, pet moderation,
adoption approval/rejection, community and lost-and-found moderation, and breed
and veterinarian catalogue administration. \\
\bottomrule
\end{tabular}
\caption{User roles and permissions.}
\label{tab:roles}
\end{table}

The role is stored in the \texttt{User} document and is never accepted from the
client. The \texttt{protect} middleware authenticates the request from the bearer
token, and the \texttt{adminOnly} middleware verifies
\texttt{req.user.role === "admin"} before any administrative route executes
(Section~\ref{sec:security}).

\section{Functional Requirements}

The functional requirements below are grouped by subsystem. Each requirement is
directly supported by a verified controller and route in the backend.

\begin{center}
\begin{longtable}{@{}p{1.1cm}p{3.4cm}p{10.0cm}@{}}
\caption{Functional requirements.}
\label{tab:fr}\\
\toprule
\textbf{ID} & \textbf{Subsystem} & \textbf{Requirement} \\
\midrule
FR-01 & Authentication & The system shall allow a user to register with name,
e-mail, and password; store the password as a bcrypt hash; and require e-mail
verification before the account can log in. \\
FR-02 & Authentication & The system shall issue a signed JWT on successful
login and require it as a bearer token on protected routes. \\
FR-03 & Authentication & The system shall support password reset through a
time-limited, single-use token sent by e-mail, and password change with the
current password. \\
FR-04 & Profile & The user shall be able to read and update own profile
(name, phone, address, city, avatar). \\
FR-05 & Pets & The user shall be able to create, list, view, update, and delete
own pets, and only own pets for protected operations. \\
FR-06 & Pets & On creation, the system shall generate a unique digital pet
identifier and a QR code encoding public pet information. \\
FR-07 & Pets & Public listing shall support filtering by species, breed,
gender, and status, full-text search, and sorting. \\
FR-08 & Health & The user shall be able to create, read, update, and delete
health records for own pets. \\
FR-09 & Vaccinations & The user shall be able to create, read, update, delete,
and list upcoming vaccinations for own pets. \\
FR-10 & Appointments & The user shall be able to book, list, update, and cancel
appointments for own pets with an active veterinarian. \\
FR-11 & Appointments & The system shall reject a double booking when the same
veterinarian already has a pending or confirmed appointment at the same
date and time. \\
FR-12 & Appointments & Booking an appointment shall automatically create a
notification and a non-duplicate reminder for the owner. \\
FR-13 & Reminders & The user shall be able to create, list, update, complete,
and delete reminders; reminders are scoped to the owner. \\
FR-14 & Notifications & The user shall be able to list, mark as read, mark all
as read, and delete own notifications, and read the unread count. \\
FR-15 & Adoption & A user shall be able to apply to adopt an available pet with
contact and motivation details, and a duplicate pending request for the same pet
shall be rejected. \\
FR-16 & Adoption & An administrator shall be able to approve or reject an
adoption request; approval automatically marks the pet as adopted and sets its
status to \texttt{adopted}. \\
FR-17 & Adoption & Any status change of an adoption request shall create a
notification for the applicant. \\
FR-18 & Favorites & The user shall be able to add and remove favorite pets; the
same user and pet pair must not be duplicated. \\
FR-19 & Community & Any visitor shall be able to read active posts; a signed-in
user shall be able to create, update, and delete own posts, and like and comment
on posts. \\
FR-20 & Lost and Found & A signed-in user shall be able to create, update, and
delete lost/found reports, and visitors shall be able to view active reports. \\
FR-21 & AI assistant & A signed-in user shall be able to ask the PetGPT
assistant questions; the answer shall consider the user's pets and fall back to
safe generic guidance when the AI service is unavailable. \\
FR-22 & AI advice & A signed-in user shall be able to request rule-based advice
for an own pet based on its vaccination status, age, and weight. \\
FR-23 & Admin & The administrator shall be able to view dashboard statistics,
list and block users, delete pets, moderate community and lost-and-found content,
approve or reject adoptions, and manage breeds and veterinarians. \\
FR-24 & E-mail & The system shall send verification, resend-verification, and
password-reset e-mails through Nodemailer, and construct links using the origin
of the request. \\
\bottomrule
\end{longtable}
\end{center}

The functional requirements table above is enforced through the route handlers
and controllers described in Chapters~5--9.

\section{Non-Functional Requirements}

\begingroup
\renewcommand{\arraystretch}{1.25}
\begin{table}[htbp]
\centering
\begin{tabular}{@{}p{3.0cm}p{11.5cm}@{}}
\toprule
\textbf{Category} & \textbf{Requirement} \\
\midrule
Security & Passwords stored as bcrypt hashes (cost factor 12); tokens (e-mail
verification, password reset) are single-use, time-limited, and stored only in
SHA-256 hashed form. \\
Access control & All protected operations are scoped to the authenticated user;
administrative operations are additionally gated by role. \\
Reliability & The system must not rely on the availability of external services:
AI and image-hosting failures degrade gracefully to stored, safe alternatives. \\
Validation & The backend validates required fields, e-mail format, password
length, resource ownership, and invalid IDs and returns appropriate HTTP errors
(400, 401, 403, 404, 500). \\
Performance & JSON responses are compressed with \texttt{compression};
appointment double-booking is prevented by a database-backed conflict query with
compound indexing on the veterinarian, date, and time. \\
Usability & The frontend is responsive and supports both light and dark theme
variants; migrated React pages are verified at 390, 768, and 1440 pixels. \\
Maintainability & The backend follows a layered structure
(routes, controllers, models, middleware, config) and a phase-based development
roadmap with verification before each commit. \\
\bottomrule
\end{tabular}
\caption{Non-functional requirements.}
\label{tab:nfr}
\end{table}
\endgroup

\section{System Scope}

The implemented scope of the system is the authenticated pet-care and adoption
workflow described by the functional requirements above, with administration
capabilities. Outside the current scope are features mentioned in configuration
or documentation but not present in the implementation, including:

\begin{itemize}
    \item the OpenAI adapter advertised in \path{backend/.env.example}
    (\path{PETGPT\_PROVIDER=openai}); the implemented assistant uses only the
    Google Gemini provider;
    \item background job processing: the \texttt{node-cron} and
    \texttt{socket.io} packages are declared dependencies but are not used by any
    application code, so no scheduled jobs, queues, or workers exist;
    \item real-time push notifications (notifications are database-bound and
    read through the API);
    \item the remaining roadmap phases 23--29 (API integration layer,
    authentication consolidation, UI completion, visual and functional
    regression, and Docker/Nginx integration, followed by
    client consolidation).
\end{itemize}

\section{Use Cases}

The use cases consolidate the functional requirements. The catalogue is organised
by actor.

\subsection{Visitor and Registered User Use Cases}

\begin{center}
\begin{longtable}{@{}p{2.0cm}p{3.6cm}p{8.9cm}@{}}
\caption{Visitor and registered-user use cases.}
\label{tab:uc-user}\\
\toprule
\textbf{Use case} & \textbf{Requirement} & \textbf{Description and key flow} \\
\midrule
UC-01 Register & FR-01 & Submit name, e-mail, password; backend normalizes the
e-mail, validates format, hashes the password, stores a hashed verification
token, and e-mails a verification link. \\
UC-02 Verify e-mail & FR-01 & Following a clicked link, the token is hashed and
matched; the account becomes verified and the token is cleared. \\
UC-03 Login & FR-02, FR-03 & Credentials are compared against the stored hash;
blocked and unverified accounts are rejected; a JWT is returned. \\
UC-04 Reset password & FR-03 & A link with a single-use token is e-mailed; the
new password is hashed and the token cleared. \\
UC-05 Manage profile & FR-04 & Read or update own profile fields. \\
UC-06 Manage pets & FR-05, FR-06, FR-07 & Create pets (QR identity generated),
list own pets, view public pets, update or delete only own pets. \\
UC-07 Manage health & FR-08, FR-09 & Create and manage health records and
vaccinations for own pets. \\
UC-08 Book appointments & FR-10--FR-12 & Book with slot conflict check; automatic
notification and reminder; reschedule or cancel keeps the reminder in sync. \\
UC-09 Manage reminders & FR-13 & Create, complete, and manage own care reminders. \\
UC-10 Manage notifications & FR-14 & View, mark read, and manage own
notifications. \\
UC-11 Apply for adoption & FR-15--FR-17 & Submit an adoption application; track
requests; receive notifications on approval or rejection. \\
UC-12 Favourites & FR-18 & Toggle pets in a personal favourites list. \\
UC-13 Community & FR-19 & Read, create, update, delete posts; like and comment. \\
UC-14 Lost and found & FR-20 & Create, update, resolve, and delete own reports;
view active reports. \\
UC-15 AI assistant & FR-21, FR-22 & Ask the pet-care assistant; request advice
for an own pet. \\
\bottomrule
\end{longtable}
\end{center}

\subsection{Administrator Use Cases}

\begin{center}
\begin{longtable}{@{}p{2.4cm}p{3.2cm}p{8.9cm}@{}}
\caption{Administrator use cases.}
\label{tab:uc-admin}\\
\toprule
\textbf{Use case} & \textbf{Requirement} & \textbf{Description} \\
\midrule
UC-16 Dashboard & FR-23 & Aggregate user, pet, adoption, lost-and-found, and
community statistics. \\
UC-17 User management & FR-23 & List users, view recent users, block/
unblock, and delete users. \\
UC-18 Pet moderation & FR-23 & List all pets with owner information; delete pet
records. \\
UC-19 Adoption moderation & FR-16, FR-17 & Approve or reject adoption
applications; approval flips the pet to \texttt{adopted}. \\
UC-20 Content moderation & FR-23 & Activate/deactivate and delete community
posts; resolve and delete lost-and-found reports. \\
UC-21 Catalogue management & FR-23 & Create, update, and delete breeds and
veterinarians. \\
\bottomrule
\end{longtable}
\end{center}

\section{Actors in the Implemented Workflow}

The system context and actor interactions are summarized in the context diagram
of Chapter~1. Community, lost-and-found, and veterinarian reading endpoints are
public, which matches the two-level access model: authenticated creation and
ownership-scoped management for content, and public read access to active
community, lost-and-found, and veterinary records.